% Page 077

\seite{77}

\margmark{27. April.}%
Gestern fand die Reichspräsidentenwahl —\ob
2.\ Wahlgang — statt. Über den ersten Wahlgang\ob
waren in unserer Gemeinde 140 Wahlberechtigte mit\ob
Stimmrecht. Gewählt haben 108 oder 77,0% (am 29.\ 3.: 109\,=\,77,9%).\ob
Ungültig 4, gültige Stimmen 104.\ Der Links\ohb
blockkandidat Marx erhielt 85 Stimmen (89), der\ob
Rechtsblockkandidat v.\ Hindenburg 17 (5). Rechnet\ob
man die Stimmen Marx u.\ Braun (Sp.) {[des]} 1. {[Wahlgangs]}\ob
zusammen\,=\,89 + 4\,=\,93, so ergibt der Linksblock\ob
gegenüber des 1.\ Wahlganges 93 – 86\,=\,7 Stimmen,\ob
während der Rechtsblock 13 Stimmen gewonnen.\ob
Dazu zu bemerken ist, daß der Kommunistische Kan\ohb
didat Thälmann im gestrigen Wahlgang nur 1\ob
Stimme erhielt. Eigentlich dürften für v.\ Hin\ohb
denburg viele Stimmen hinzugekommen sein.

\margmark{Mai}%
Das amtliche Resultat des 2.\ Wahlganges der Reichs\ohb
präsidentenwahl v.\ 26.\ 4. 1925 ist folgendes:

\pend
\begin{verse}
Gewählt haben:\\
3 0 5 8 1 4 1 5\\
davon sind gültige Stimmen: 3 0 3 5 1 9 4 8\\
Zahl der ungültigen St.:      2 1 6 0 5 1
\end{verse}
\pstart

Rechtsblockkandidat v.\ Hindenburg: 1 4 6 5 5 7 6 6\ob
Linksblockkandidat Marx           : 1 3 7 5 1 6 1 5\ob
Kommunist. Kand. Thaelmann        :   1 9 3 1 1 5 1

Somit ist v.\ Hindenburg zum Reichspräsidenten\ob
gewählt.

\margmark{1. Juni.}%
Freudig zu begrüßen ist es, daß der Weg durch den\ob
Acker vollständig neu gemacht worden ist. Auch die\ob
Langstraße direkt durchs Dorf ward in Kriegszeiten\ob
eingearbeitet, da die Vorbereitungen in\ob
\unsicher bereits vollendet sind.

\margmark{15. Juni.}%
8 Tage Pfingstferien wurden an die Feiertage an\ohb
gelegt. Da Fronleichnam darin fällt, somit wäre\ob
17.\ Juni der Aufpfingsttag, 8.\ Juli.

\margmark{16. Juni.}%
Aus Anlaß der heutigen Volks-, Betriebs- und Berufs\ohb
zählung fällt der Unterricht aus.
\pend
